% ============================================================
% SECTION 0 — The Puzzle & the Interpretive Key (Introduction)
% ============================================================
\section{The Puzzle}

\begin{frame}{Why Do the Imams Look So Different?}
\bigquote{We are facing a matter that has several sides and colors, which neither hearts can bear nor intelligence fathom.}{'Ali, on being asked to accept the caliphate}
\begin{itemize}[itemsep=4pt]
  \item al-Hassan $\rightarrow$ negotiates peace with Mu'awiyah
  \item al-Husayn $\rightarrow$ revolts against Yazid to the death
  \item Zayn al-'Abidin $\rightarrow$ worship, service, decades of mourning
  \item al-Sadiq $\rightarrow$ scholarship, refuses an offered throne
  \item Musa al-Kazim $\rightarrow$ endures four years of imprisonment
  \item al-Rida $\rightarrow$ accepts the crown-princeship
  \item al-'Askari $\rightarrow$ leads under total state surveillance
\end{itemize}
\delivernote{Let the room actually sit in the contradiction for a moment before moving on: al-Husayn dies rather than submit; al-Hassan negotiates with the very same enemy family. Ask: can both be right?}
\end{frame}

\begin{frame}{The Dilemma Taqiyyah Creates}
\bigquote{Leave me and seek someone else.}{'Ali, Nahj al-Balaghah, Sermon 91}
\textbf{Taqiyyah} --- concealing one's true belief or intent under pressure --- is closely tied to Shi'ism, and several Imams used it.
\begin{itemize}[itemsep=4pt]
  \item If taqiyyah is legitimate $\rightarrow$ why did al-Husayn refuse it and die?
  \item If taqiyyah is illegitimate $\rightarrow$ why did 'Ali, al-Hassan, al-Sajjad, and al-Sadiq all practice and recommend it?
\end{itemize}
\examplenote{Fourteen doctors}{Shi'ism accepts 14 infallibles as authoritative --- more hadiths than even the Sunni \emph{Sihah al-Sittah} (al-Kafi alone: 16{,}000+). Richness without careful research risks chaos: ``once a patient has too many doctors, there is no hope for recovery.''}
\end{frame}

\begin{frame}{The Interpretive Key: al-jam' al-'urfi}
\textbf{al-jam' al-'urfi} (``conventional reconciliation''): apparent contradictions dissolve once you separate a fixed \emph{essence} from its circumstance-driven \emph{expression}.

\vspace{0.3cm}
\begin{center}
\begin{tikzpicture}[
  every node/.style={font=\small,align=center},
  ess/.style={rectangle,rounded corners,draw=deepteal,thick,fill=lightbg,text width=3.6cm,minimum height=1.1cm},
  exp/.style={rectangle,rounded corners,draw=gold,thick,fill=white,text width=2.6cm,minimum height=1cm}
]
\node[ess] (ess) {\color{deepteal}\bfseries Essence (fixed)\\ sympathy, compassion, equality, justice, fairness};
\node[exp, below left=1.1cm and 2.3cm of ess] (e1) {'Ali: poverty, patched clothes};
\node[exp, below=1.1cm of ess] (e2) {al-Sadiq: fine clothes \emph{and} famine bread};
\node[exp, below right=1.1cm and 2.3cm of ess] (e3) {'Ali: dye tradition suspended in war};
\draw[-{Latex}, thick, deepteal] (ess) -- (e1);
\draw[-{Latex}, thick, deepteal] (ess) -- (e2);
\draw[-{Latex}, thick, deepteal] (ess) -- (e3);
\end{tikzpicture}
\end{center}
\delivernote{Ask: what stays fixed across every branch, and what is free to change? (Answer: essence fixed; outward expression adapts to what the moment requires.)}
\end{frame}

\begin{frame}{Two Proofs 'Ali Gives Us Himself}
\begin{itemize}[itemsep=6pt]
  \item \textbf{al-Sadiq's clothing:} accused of departing from the Prophet's simplicity, he answers Islam commands \emph{values}, not a wardrobe --- then proves it by selling his own stocked wheat during a famine.
  \item \textbf{'Ali's beard-dyeing hadith:} he narrated ``color the white hairs... do not resemble the Jews'' --- but never dyed his own beard. His own explanation: it was a wartime tactic (hide who was old, who was young), not a permanent command.
\end{itemize}
\bigquote{If your father was lacking common sense, does that mean you should be foolish like he was?}{Imam al-Sadiq, to a man living in a cramped inherited house}
\end{frame}

\qaframe{The Puzzle \& Interpretive Key}{%
  \qa{What is al-jam' al-'urfi, and why does the book need it before Chapter 1 even begins?}
     {It is the rule that apparent contradictions between the Imams dissolve once essence is separated from expression. Without it, every later chapter's central puzzle would look like the Imams simply disagreeing with each other.}
  \qa{Why is 'Ali narrating a hadith he never personally practiced the stronger of the two proofs?}
     {Because the same person is both the narrator of the instruction and the one who visibly didn't follow it -- proving the gap was principled and understood, not hypocrisy or forgetfulness.}
  \qa{The ``too many doctors'' line criticizes abundance, yet fourteen Imams is also framed as a strength. Can both be true?}
     {Yes -- it is a double-edged strength: more guidance across more centuries, but also more surface area for forged attribution. The reconciliation method this Introduction builds exists to manage that second edge.}
}
